% from ega4-2-fr.tex
\oldpage[IV-2]{183}

\begin{definition}[7.1.1]
\label{IV.2.s7.1.1-fr}
On dit qu'un anneau local noethérien $A$ est \emph{formellement
équidimensionnel} si son complété $\widehat A$ est équidimensionnel (0,
16.1.4).
\end{definition}

\begin{proposition}[7.1.2]
\label{IV.2.s7.1.2-fr}
Soient $A$ un anneau local noethérien, $\mathfrak p_i$ ($1\leqslant i\leqslant
n$) ses idéaux premiers minimaux. Pour que $A$ soit formellement
équidimensionnel, il faut et il suffit que les $A/\mathfrak p_i$ le soient et
que $A$ soit équidimensionnel (autrement dit que les $A/\mathfrak p_i$ aient
même dimension).
\end{proposition}

\begin{proof}
En effet, cela résulte du fait que pour tout idéal premier $\mathfrak p'$ de
$\widehat A$, $\mathfrak p'\cap A$ contient l'un des $\mathfrak p_i$, donc
$\mathfrak p'$ contient l'un des $\mathfrak p_i\widehat A$, et
$\widehat A/\mathfrak p_i\widehat A=(A/\mathfrak p_i)\otimes_A\widehat A$
est le complété de l'anneau local $A/\mathfrak p_i$ (0$_{\mathrm I}$,
7.3.3), donc a même dimension. Toute chaîne maximale d'idéaux premiers de
$\widehat A$ s'identifie donc canoniquement à une chaîne maximale d'idéaux
premiers de l'un des $\widehat A/\mathfrak p_i\widehat A$, et réciproquement,
d'où aussitôt la conclusion.
\end{proof}

\begin{proposition}[7.1.3]
\label{IV.2.s7.1.3-fr}
Soient $A$, $A'$ deux anneaux locaux noethériens, $\mathfrak m$ l'idéal
maximal de $A$, $\varphi:A\to A'$ un homomorphisme local; on suppose que $A'$
est un $A$-module plat, et que $A'$ soit équidimensionnel et caténaire. Alors:
\begin{enumerate}
\item[(i)] $A$ est équidimensionnel et caténaire.
\item[(ii)] Supposons en outre que $\mathfrak mA'$ soit un idéal de définition de
$A'$. Alors, si $\mathfrak a$ est un idéal de $A$, pour que
$A'/\mathfrak aA'$ soit équidimensionnel, il faut et il suffit que
$A/\mathfrak aA$ le soit. En particulier, pour tout idéal premier
$\mathfrak p$ de $A$, $A'/\mathfrak pA'$ est équidimensionnel.
\end{enumerate}
\end{proposition}

\begin{proof}
Prouvons d'abord la seconde assertion de (ii). Posons
$X=\operatorname{Spec}(A)$, $X'=\operatorname{Spec}(A')$,
$Y=V(\mathfrak p)=\operatorname{Spec}(A/\mathfrak p)$,
$Y'=V(\mathfrak pA')=\operatorname{Spec}(A'/\mathfrak pA')$; si
$f:X'\to X$ est le morphisme correspondant à $\varphi$, $Y'$ est aussi le
sous-préschéma fermé $f^{-1}(Y)$ de $X'$. Comme $\mathfrak mA'$ est un idéal
de définition de $A'$, on a
\begin{equation}
\dim(X)=\dim(X')
\tag{7.1.3.1}
\label{IV.2.eq.7.1.3.1-fr}
\end{equation}
en vertu de (6.1.3); de même $\mathfrak mA'/\mathfrak pA'$ est un idéal de
définition de $A'/\mathfrak pA'$ et $Y'$ est plat sur $Y$ (2.1.4), donc
(6.1.3)
\begin{equation}
\dim(Y)=\dim(Y').
\tag{7.1.3.2}
\label{IV.2.eq.7.1.3.2-fr}
\end{equation}
Soient $y$ le point générique de $Y$, $Y'_i$ ($1\leqslant i\leqslant r$) les
composantes irréductibles de $Y'$, $y'_i$ le point générique de $Y'_i$
($1\leqslant i\leqslant r$); on a $f(y'_i)=y$ pour tout $i$ (2.3.4).
D'autre part, $y'_i$ est point maximal de la fibre $f^{-1}(y)$ (0$_{\mathrm
I}$, 2.1.8), donc
$\mathscr O_{X',y'_i}\otimes_{\mathscr O_{X,y}}k(y)=
\mathscr O_{X',y'_i}/\mathfrak m_y\mathscr O_{X',y'_i}$ est de dimension
0, autrement dit $\mathfrak m_y\mathscr O_{X',y'_i}$ est un idéal de
définition de $\mathscr O_{X',y'_i}$; on peut de nouveau appliquer (6.1.3),
qui donne
\begin{equation}
\dim(\mathscr O_{X',y'_i})=\dim(\mathscr O_{X,y}),
\tag{7.1.3.3}
\label{IV.2.eq.7.1.3.3-fr}
\end{equation}
c'est-à-dire (5.1.2)
\begin{equation}
\operatorname{codim}(Y'_i,X')=\operatorname{codim}(Y,X)
\tag{7.1.3.4}
\label{IV.2.eq.7.1.3.4-fr}
\end{equation}
pour tout $i$. Comme $A'$ est équidimensionnel et caténaire, on a, en vertu
de (0, 14.3.5)
\begin{equation}
\dim(X')=\dim(Y'_i)+\operatorname{codim}(Y'_i,X')
\tag{7.1.3.5}
\label{IV.2.eq.7.1.3.5-fr}
\end{equation}

